% New English translation from checked cote 25, site 102 (PDF page 103).
% AI-translated by GPT-6 Astra (OpenAI, ChatGPT Pro mode).
\scanpage{25/102}
\mspage{5}
% NOTE: the reviewed source records three long diagonal pen strokes; the text remains live beneath them.
of Picard, obtained from the section of the latter defined by $\underline{O}_P(1)$, is an
\uline{isomorphism}.

Let P be a projective bundle over a field k. An invertible Module
$\underline{L}$ on P is said to be \uline{of degree} n if $\underline{L}$ is isomorphic to $\underline{O}_P(n)$; when
dim P$\hand{\geq}$1, this determines n in terms of $\underline{L}$, but if dim P $\hand{\leq}$ D
(i.e. p is empty, or reduced to a point) then $\underline{L}$ is degree n for every n.
When P is a projective bundle over an arbitrary base, an invertible
Module $\underline{L}$ on P is said to be of degree n if it is of degree n on every fibr
% NOTE: the end of the preceding source line is cut off at the right edge of the scan.
This determines n, except if all the fibers of P/S are of dim$\hand{\leq}$0, then
$\underline{L}$ is of degree n for every n. To say that $\underline{L}$ is of degree n also means,
by 1.1. and 1.2., that \var{\inserted[typed]{the class}}{the class}{insertion}{ega25_25_102-v0002} \var{\xout{the image}}{}{typewriter strike-out}{ega25_25_102-v0003} of $\underline{L}$ in Pic(P) lies in
the image of \var{\inserted{$\mathrm{Pic}\,S\times\{n\}$ \unsure{by|by}}}{$\mathrm{Pic}\,S\times\{n\}$ \unsure{by|by}}{insertion}{ega25_25_102-v0004} the homomorphism (x) Pic(S)x$\underline{Z}$ $\hand{\longrightarrow}$Pic(P) considered above,
i.e. that $\underline{L}$ is isomorphic to a Module of the form $f^*(\underline{M})\unsure{(|\otimes}n)$, where $\underline{M}$ is
invertible on S. Moreover, when the fibers of P are nonempty
i.e..$\underline{E}$ everywhere of rank $\hand{\geq}$1, then $\underline{M}$ is determined up to unique isomorphism\var{\margin{\ensuremath{/}}}{\placedmargin{\ensuremath{/}}}{marginal note placed}{ega25_25_102-v0005}
\var{\margin{\textbar}}{\placedmargin{\textbar}}{marginal note placed}{ega25_25_102-v0006}in terms of $\underline{L}$, as follows again from 1.1. and 1.2. (On this
subject, I realize that 1.1. should have been stated thus, without a hypothesis
on the rank of $\underline{E}$: every invertible $\underline{L}$ on S can be put in the form
indicated in the statement; when the fibers of P are nonempty i.e. the
rank of $\underline{E}$ everywhere $\hand{\neq}$ 0, then $\underline{M}$ is determined up to unique isom by
$\underline{L}$; when the fibers of P are all of dim $\hand{\geq}$1, \var{\inserted{i.e. the rank of $\underline{E}$ everywhere $\geq 2$,}}{i.e. the rank of $\underline{E}$ everywhere $\geq 2$,}{insertion}{ega25_25_102-v0007} then the partition
of S is also uniquely determined by the choice of $\underline{L}$. In this
way, Remark 1. 2. is eliminated and passes into the proof).\var{\margin{\ensuremath{\swarrow}}}{\placedmargin{\ensuremath{\swarrow}}}{marginal note placed}{ega25_25_102-v0008}

Let P' = P($\underline{E}$') be a second projective bundle; then the determination
of Pic(P) allows us in principle to determine the S-morphisms
g:P$\hand{\to}$P', since these are defined by an invertible Module $\underline{L}$ \var{\inserted{($\simeq g^*\mathcal{O}_{P'}(1)$)}}{($\simeq g^*\mathcal{O}_{P'}(1)$)}{insertion}{ega25_25_102-v0009} on P\var{\margin{\ensuremath{/}}}{\placedmargin{\ensuremath{/}}}{marginal note placed}{ega25_25_102-v0010}
\tnote{25-102-1}{The source's capital D in the dimension bound, lowercase p in its parenthesis, the clipped fragment fibr, and the statement invertible on S in the proposed reformulation are retained. The abbreviation alternatives pr or par both render by; their uncertainty remains marked. The removed stray t before termes has no English spelling effect.}

\var{\margin{102}}{\placedmargin{102}}{marginal note placed}{ega25_25_102-v0011}
